\subsection{Pure extended-real antitone limits and genuine rational approximation}

\begin{lemma}[Actual infima and antitone extended-real arithmetic]
\label{analysis:actual-ereal-antitone-infima}
\leanok
\lean{EnvelopeProperty.iInf_eq_of_tendsto_of_antitone,
EnvelopeProperty.iInf_comp_of_antitone,
EnvelopeProperty.iInf_max_of_antitone,
EnvelopeProperty.iInf_add_of_antitone,
EnvelopeProperty.iInf_add_coe_of_antitone,
EnvelopeProperty.iInf_coe_mul_of_antitone}
The actual infimum of a decreasing extended-real sequence is its limit. A
monotone map continuous at that original infimum commutes with it. Finite
maxima commute with the actual infima of two antitone sequences. Their sums
do so when each sequence has a first term different from positive infinity;
adding any finite real constant also commutes. Multiplication by a strictly
positive finite real commutes with any such antitone infimum. All statements
use the ORIGINAL extended-real operations, including their actual infinity
conventions; no convergence or arithmetic certificate is supplied.
\end{lemma}
\begin{proof}
The order-topology antitone convergence theorem gives convergence to the
actual infimum. If an independently given limit exists, uniqueness in the
Hausdorff extended-real topology identifies it with this same infimum.
Continuous composition and monotonicity then identify the actual image
infimum. For maxima use genuine continuity of maximum and the actual product
of the two convergent sequences.

For sums, the original infima are bounded above by their respective first
terms different from positive infinity and hence cannot be positive infinity. At this actual pair,
extended-real addition is continuous: neither mixed positive/negative infinity
case occurs. Compose that genuine continuity with the original two sequence
limits; the sums remain antitone, so uniqueness gives exactly their original
infimum. A constant finite real is a particular constant antitone sequence.
For multiplication by a fixed positive finite real, the actual extended-real
map is monotone and continuous, including at both infinities. Apply the same
continuous-composition argument. This retains negative-infinite limits and
never assumes uniform convergence.
\end{proof}

\begin{lemma}[True negative-infinite examples and rational approximation of every real]
\label{analysis:actual-rational-antitone-real-approximation}
\leanok
\lean{EnvelopeProperty.iInf_coe_sub_natCast,
EnvelopeProperty.iInf_neg_natCast,
EnvelopeProperty.coe_mul_ne_bot,
EnvelopeProperty.exists_antitone_rat_iInf_eq}
For every real $c$, the original infimum of $c-j$ over natural $j$ is negative
infinity, as is the infimum of $-j$. A positive finite real multiple of an
extended real different from negative infinity is still different from
negative infinity. Every real $c$ admits an ACTUAL antitone rational sequence
$q_j$ with $c<q_j$ and actual extended-real infimum equal to $c$.
\end{lemma}
\begin{proof}
For every finite real threshold choose a genuine natural integer large enough
that $c-j$ lies below it. The actual characterization of infimum at negative
infinity proves the first assertion and its $c=0$ specialization. The
multiplication assertion follows by the three actual extended-real cases:
negative infinity is excluded, a finite real product remains finite, and a
strictly positive finite real times positive infinity is positive infinity.

For each natural $j$, rational density chooses an ORIGINAL rational $q_j$ in
\[
 c+\frac1{j+2}<q_j<c+\frac1{j+1}.
\]
The adjacent open intervals give $q_{j+1}<q_j$, hence the actual sequence is
antitone, and every term lies above $c$. This gives the lower bound $c$ for
its actual infimum. Given any real threshold $b>c$, take a true natural $j$
with $1/(j+1)<b-c$. The displayed upper bound gives $q_j<b$, so the actual
infimum is at most every such threshold and therefore at most $c$. The
positive-infinite threshold is immediate from any finite first term, and the
negative-infinite threshold cannot lie above $c$. Thus the original infimum
is exactly the original $c$.
\end{proof}
